% GENERATED FILE. Edit canonical lessons, then run npm run generate:books.
% canonical-source-hash: fnv1a64:36bd326aeffe0f99
% canonical-lessons: PA-C259-nambar, PA-C259-ghanti, PA-C259-vajjna, PA-C259-avaz, PA-C259-milauna, PA-R259-review-253-259

\chapter{When the Phone Rings}
\label{ch:pa-a2-when-the-phone-rings}
\hlchaptermodality{259}

\begin{chapteropening}
\textbf{By the end of this chapter:} \emph{I can give a phone number and talk about a call.}

Complete the last lesson of chapter 259: I can give a phone number and talk about a call.
\end{chapteropening}

\section[\texorpdfstring{\pa{ਨੰਬਰ} (nambar)}{ਨੰਬਰ (nambar)}]{\pa{ਨੰਬਰ} (nambar) — a number; a phone number}

\label{lesson:PA-C259-nambar}



\begin{quote}
Before the new one: say the Punjabi for a banknote, then the Punjabi for to earn.
\end{quote}

\subsection*{You'll want to know: \pa{ਨੰਬਰ}}
\textbf{\pa{ਨੰਬਰ}} — \emph{nambar} — \textquotedblleft{}a number; a phone number\textquotedblright{}.

\begin{quote}
\pa{ਤੁਹਾਡਾ} \pa{ਫ਼ੋਨ} \pa{ਨੰਬਰ} \pa{ਕੀ} \pa{ਹੈ}?
\end{quote}

\begin{quote}
\emph{tuhāḍā fon nambar kī hai?}
\end{quote}

\begin{quote}
What is your phone number?
\end{quote}

A masculine noun.

\begin{cousinweb}
\textbf{\pa{ਨੰਬਰ}} (\emph{nambar}) is borrowed from English \emph{number}.
\end{cousinweb}

\begin{grammarlens}[title={What you've built}]
One more word for reading notices and messages.
\end{grammarlens}

\subsection*{Your turn}
\begin{itemize}
\raggedright
  \item \emph{Say it:} \emph{nambar}
  \item \emph{Say it:} \emph{nambar}, once more
  \item \emph{Say it:} \emph{tuhāḍā fon nambar kī hai?} — What is your phone number?
  \item \emph{Recall:} say the Punjabi for a banknote, then the Punjabi for to earn, then say \emph{nambar} again
\end{itemize}

\begin{tcolorbox}[breakable,colback=teal!4,colframe=teal!35!black,title={Before you move on}]
What does \pa{ਨੰਬਰ} mean? (\textquotedblleft{}A number; a phone number\textquotedblright{}.) Say it once more.
\end{tcolorbox}
\section[\texorpdfstring{\pa{ਘੰਟੀ} (ghaṇṭī)}{ਘੰਟੀ (ghaṇṭī)}]{\pa{ਘੰਟੀ} (ghaṇṭī) — a bell; a ring}

\label{lesson:PA-C259-ghanti}



\begin{quote}
Before the new one: say the Punjabi for to earn, then the Punjabi for a number.
\end{quote}

\subsection*{You'll want to know: \pa{ਘੰਟੀ}}
\textbf{\pa{ਘੰਟੀ}} — \emph{ghaṇṭī} — \textquotedblleft{}a bell; a ring\textquotedblright{}.

\begin{quote}
\pa{ਫ਼ੋਨ} \pa{ਦੀ} \pa{ਘੰਟੀ} \pa{ਬਹੁਤ} \pa{ਉੱਚੀ} \pa{ਹੈ।}
\end{quote}

\begin{quote}
\emph{fon dī ghaṇṭī bahut uchchī hai.}
\end{quote}

\begin{quote}
The phone's ring is very loud.
\end{quote}

A feminine noun.

\begin{cousinweb}
\textbf{\pa{ਘੰਟੀ}} (\emph{ghaṇṭī}) is inherited from Sanskrit \emph{ghaṇṭikā}, \textquotedblleft{}a small bell\textquotedblright{}, beside \emph{ghaṇṭā}, \textquotedblleft{}bell\textquotedblright{}; CDIAL lists Punjabi \emph{ghaṇḍā}, \textquotedblleft{}a bell used to attract deer\textquotedblright{}, under it (CDIAL 4421).
\end{cousinweb}

\begin{grammarlens}[title={What you've built}]
One more word for reading notices and messages.
\end{grammarlens}

\subsection*{Your turn}
\begin{itemize}
\raggedright
  \item \emph{Say it:} \emph{ghaṇṭī}
  \item \emph{Say it:} \emph{ghaṇṭī}, once more
  \item \emph{Say it:} \emph{fon dī ghaṇṭī bahut uchchī hai.} — The phone's ring is very loud.
  \item \emph{Recall:} say the Punjabi for to earn, then the Punjabi for a number, then say \emph{ghaṇṭī} again
  \item \emph{Write it:} copy \textbf{\pa{ਘੰਟੀ}} once, checking the whole word
\end{itemize}

\begin{tcolorbox}[breakable,colback=teal!4,colframe=teal!35!black,title={Before you move on}]
What does \pa{ਘੰਟੀ} mean? (\textquotedblleft{}A bell; a ring\textquotedblright{}.) Say it once more.
\end{tcolorbox}
\section[\texorpdfstring{\pa{ਵੱਜਣਾ} (vajjṇā)}{ਵੱਜਣਾ (vajjṇā)}]{\pa{ਵੱਜਣਾ} (vajjṇā) — to ring; to sound}

\label{lesson:PA-C259-vajjna}



\begin{quote}
Before the new one: say the Punjabi for a number, then the Punjabi for a bell.
\end{quote}

\subsection*{You'll want to know: \pa{ਵੱਜਣਾ}}
\textbf{\pa{ਵੱਜਣਾ}} — \emph{vajjṇā} — \textquotedblleft{}to ring; to sound\textquotedblright{}.

\begin{quote}
\pa{ਰਾਤ} \pa{ਨੂੰ} \pa{ਫ਼ੋਨ} \pa{ਨਹੀਂ} \pa{ਵੱਜਦਾ} \pa{ਹੈ।}
\end{quote}

\begin{quote}
\emph{rāt nū̃ fon nahīṁ vajjdā hai.}
\end{quote}

\begin{quote}
The phone does not ring at night.
\end{quote}

Past (it …): \pa{ਉਹ} \pa{ਘੰਟੀ} \pa{ਵੱਜੀ}, \emph{oh ghaṇṭī vajjī}.

\begin{cousinweb}
\textbf{\pa{ਵੱਜਣਾ}} (\emph{vajjṇā}) is inherited from Sanskrit \emph{vādyate}, \textquotedblleft{}is made to sound\textquotedblright{}; CDIAL lists Punjabi \emph{vajjṇā}, \textquotedblleft{}to sound (of bell, clock)\textquotedblright{}, under it (CDIAL 11513).
\end{cousinweb}

\begin{grammarlens}[title={What you've built}]
One more word for reading notices and messages.
\end{grammarlens}

\subsection*{Your turn}
\begin{itemize}
\raggedright
  \item \emph{Say it:} \emph{vajjṇā}
  \item \emph{Say it:} \emph{vajjṇā}, once more
  \item \emph{Say it:} \emph{rāt nū̃ fon nahīṁ vajjdā hai.} — The phone does not ring at night.
  \item \emph{Recall:} say the Punjabi for a number, then the Punjabi for a bell, then say \emph{vajjṇā} again
  \item \emph{Write it:} copy \textbf{\pa{ਵੱਜਣਾ}} once, checking the whole word
\end{itemize}

\begin{tcolorbox}[breakable,colback=teal!4,colframe=teal!35!black,title={Before you move on}]
What does \pa{ਵੱਜਣਾ} mean? (\textquotedblleft{}To ring; to sound\textquotedblright{}.) Say it once more.
\end{tcolorbox}
\section[\texorpdfstring{\pa{ਆਵਾਜ਼} (āvāz)}{ਆਵਾਜ਼ (āvāz)}]{\pa{ਆਵਾਜ਼} (āvāz) — a voice; a sound}

\label{lesson:PA-C259-avaz}



\begin{quote}
Before the new one: say the Punjabi for a bell, then the Punjabi for to ring.
\end{quote}

\subsection*{You'll want to know: \pa{ਆਵਾਜ਼}}
\textbf{\pa{ਆਵਾਜ਼}} — \emph{āvāz} — \textquotedblleft{}a voice; a sound\textquotedblright{}.

\begin{quote}
\pa{ਫ਼ੋਨ} \pa{ਉੱਤੇ} \pa{ਅਮਨ} \pa{ਦੀ} \pa{ਆਵਾਜ਼} \pa{ਸਾਫ਼} \pa{ਹੈ।}
\end{quote}

\begin{quote}
\emph{fon utte Aman dī āvāz sāf hai.}
\end{quote}

\begin{quote}
Aman's voice is clear on the phone.
\end{quote}

A feminine noun.

\begin{cousinweb}
\textbf{\pa{ਆਵਾਜ਼}} (\emph{āvāz}) is a loan, usually traced to Persian \emph{āvāz}, \textquotedblleft{}a voice\textquotedblright{}.
\end{cousinweb}

\begin{grammarlens}[title={What you've built}]
One more word for reading notices and messages.
\end{grammarlens}

\subsection*{Your turn}
\begin{itemize}
\raggedright
  \item \emph{Say it:} \emph{āvāz}
  \item \emph{Say it:} \emph{āvāz}, once more
  \item \emph{Say it:} \emph{fon utte Aman dī āvāz sāf hai.} — Aman's voice is clear on the phone.
  \item \emph{Recall:} say the Punjabi for a bell, then the Punjabi for to ring, then say \emph{āvāz} again
  \item \emph{Write it:} copy \textbf{\pa{ਆਵਾਜ਼}} once, checking the whole word
\end{itemize}

\begin{tcolorbox}[breakable,colback=teal!4,colframe=teal!35!black,title={Before you move on}]
What does \pa{ਆਵਾਜ਼} mean? (\textquotedblleft{}A voice; a sound\textquotedblright{}.) Say it once more.
\end{tcolorbox}
\section[\texorpdfstring{\pa{ਮਿਲਾਉਣਾ} (milāuṇā)}{ਮਿਲਾਉਣਾ (milāuṇā)}]{\pa{ਮਿਲਾਉਣਾ} (milāuṇā) — to dial a number; to put through}

\label{lesson:PA-C259-milauna}



\begin{quote}
Before the new one: say the Punjabi for to ring, then the Punjabi for a voice.
\end{quote}

\subsection*{You'll want to know: \pa{ਮਿਲਾਉਣਾ}}
\textbf{\pa{ਮਿਲਾਉਣਾ}} — \emph{milāuṇā} — \textquotedblleft{}to dial a number; to put through\textquotedblright{}.

\begin{quote}
\pa{ਮੈਂ} \pa{ਸ਼ਾਮ} \pa{ਨੂੰ} \pa{ਅਮਨ} \pa{ਦਾ} \pa{ਨੰਬਰ} \pa{ਮਿਲਾਉਂਦਾ} \pa{ਹਾਂ।}
\end{quote}

\begin{quote}
\emph{maiṁ shām nū̃ Aman dā nambar milāundā hāṁ.}
\end{quote}

\begin{quote}
In the evening I dial Aman's number.
\end{quote}

Past (I …): \pa{ਮੈਂ} \pa{ਬੈਂਕ} \pa{ਦਾ} \pa{ਨੰਬਰ} \pa{ਮਿਲਾਇਆ}, \emph{maiṁ baiṅk dā nambar milāiā}.

\begin{cousinweb}
\textbf{\pa{ਮਿਲਾਉਣਾ}} (\emph{milāuṇā}) is the causative partner of \emph{milṇā}, \textquotedblleft{}to meet\textquotedblright{}; CDIAL lists Punjabi \emph{milṇā}, \textquotedblleft{}to meet, be obtained\textquotedblright{}, under Sanskrit \emph{milati}, \textquotedblleft{}meets\textquotedblright{} (CDIAL 10133).
\end{cousinweb}

\begin{grammarlens}[title={What you've built}]
One more word for reading notices and messages.
\end{grammarlens}

\subsection*{Your turn}
\begin{itemize}
\raggedright
  \item \emph{Say it:} \emph{milāuṇā}
  \item \emph{Say it:} \emph{milāuṇā}, once more
  \item \emph{Say it:} \emph{maiṁ shām nū̃ Aman dā nambar milāundā hāṁ.} — In the evening I dial Aman's number.
  \item \emph{Recall:} say the Punjabi for to ring, then the Punjabi for a voice, then say \emph{milāuṇā} again
  \item \emph{Write it:} copy \textbf{\pa{ਮਿਲਾਉਣਾ}} once, checking the whole word
\end{itemize}

\begin{tcolorbox}[breakable,colback=teal!4,colframe=teal!35!black,title={Before you move on}]
What does \pa{ਮਿਲਾਉਣਾ} mean? (\textquotedblleft{}To dial a number; to put through\textquotedblright{}.) Say it once more.
\end{tcolorbox}
\section[(dialogue)]{Looking back: from Rules at the Gate to When the Phone Rings}

\label{lesson:PA-R259-review-253-259}



\begin{quote}
Say the Punjabi for a rule and a notice.
\end{quote}

\subsection*{Your turn}
Say these aloud:

\begin{itemize}
\raggedright
  \item a rule, a notice, a queue, to move out of the way, an advertisement
  \item a form to fill in, an application, correct, wrong, to rub out
  \item the post, an envelope, an office stamp, paper, to print
  \item a postman, a parcel, to deliver, to send off, to tear
  \item a bank account, a sum of money, in cash, to withdraw, to go up
  \item credit, a debt, a coin, a banknote, to earn
  \item a number, a bell, to ring, a voice, to dial a number
\end{itemize}

\begin{tcolorbox}[breakable,colback=teal!4,colframe=teal!35!black,title={Before you move on}]
A rule? (\textbf{\pa{ਨਿਯਮ}}, \emph{niyam}.) A notice? (\textbf{\pa{ਨੋਟਿਸ}}, \emph{noṭis}.) A queue? (\textbf{\pa{ਕਤਾਰ}}, \emph{katār}.) To move out of the way? (\textbf{\pa{ਹਟਣਾ}}, \emph{haṭṇā}.) An advertisement? (\textbf{\pa{ਇਸ਼ਤਿਹਾਰ}}, \emph{ishtihār}.) A form to fill in? (\textbf{\pa{ਫ਼ਾਰਮ}}, \emph{fāram}.) An application? (\textbf{\pa{ਅਰਜ਼ੀ}}, \emph{arzī}.) Correct? (\textbf{\pa{ਸਹੀ}}, \emph{sahī}.) Wrong? (\textbf{\pa{ਗ਼ਲਤ}}, \emph{ghalat}.) To rub out? (\textbf{\pa{ਮਿਟਾਉਣਾ}}, \emph{miṭāuṇā}.) The post? (\textbf{\pa{ਡਾਕ}}, \emph{ḍāk}.) An envelope? (\textbf{\pa{ਲਿਫ਼ਾਫ਼ਾ}}, \emph{lifāfā}.) An office stamp? (\textbf{\pa{ਮੋਹਰ}}, \emph{mohar}.) Paper? (\textbf{\pa{ਕਾਗ਼ਜ਼}}, \emph{kāghaz}.) To print? (\textbf{\pa{ਛਾਪਣਾ}}, \emph{chhāpṇā}.) A postman? (\textbf{\pa{ਡਾਕੀਆ}}, \emph{ḍākīā}.) A parcel? (\textbf{\pa{ਪਾਰਸਲ}}, \emph{pārsal}.) To deliver? (\textbf{\pa{ਪਹੁੰਚਾਉਣਾ}}, \emph{pahunchāuṇā}.) To send off? (\textbf{\pa{ਘੱਲਣਾ}}, \emph{ghallṇā}.) To tear? (\textbf{\pa{ਪਾੜਨਾ}}, \emph{pāṛnā}.) A bank account? (\textbf{\pa{ਖਾਤਾ}}, \emph{khātā}.) A sum of money? (\textbf{\pa{ਰਕਮ}}, \emph{rakam}.) In cash? (\textbf{\pa{ਨਕਦ}}, \emph{nakad}.) To withdraw? (\textbf{\pa{ਕਢਵਾਉਣਾ}}, \emph{kaḍhvāuṇā}.) To go up? (\textbf{\pa{ਵਧਣਾ}}, \emph{vadhṇā}.) Credit? (\textbf{\pa{ਉਧਾਰ}}, \emph{udhār}.) A debt? (\textbf{\pa{ਕਰਜ਼ਾ}}, \emph{karzā}.) A coin? (\textbf{\pa{ਸਿੱਕਾ}}, \emph{sikkā}.) A banknote? (\textbf{\pa{ਨੋਟ}}, \emph{noṭ}.) To earn? (\textbf{\pa{ਕਮਾਉਣਾ}}, \emph{kamāuṇā}.) A number? (\textbf{\pa{ਨੰਬਰ}}, \emph{nambar}.) A bell? (\textbf{\pa{ਘੰਟੀ}}, \emph{ghaṇṭī}.) To ring? (\textbf{\pa{ਵੱਜਣਾ}}, \emph{vajjṇā}.) A voice? (\textbf{\pa{ਆਵਾਜ਼}}, \emph{āvāz}.) To dial a number? (\textbf{\pa{ਮਿਲਾਉਣਾ}}, \emph{milāuṇā}.)
\end{tcolorbox}
